\section{Introducción}
\label{sec:introduccion}

El análisis diferencial estudia cómo las diferencias entre dos entradas se
propagan a través de una primitiva criptográfica. Dentro de este enfoque, el
número de componentes no lineales activados constituye un indicador
estructural relevante: cuanto mayor sea la cantidad de componentes activos,
mayor será la dispersión de la diferencia y, en general, más difícil resultará
construir trayectorias diferenciales de baja complejidad
\cite{biham1991differential}.

Keccak, base del estándar SHA-3, organiza su estado como un arreglo
tridimensional de bits y aplica en cada ronda las transformaciones
$\theta$, $\rho$, $\pi$, $\chi$ e $\iota$
\cite{nist_fips202,bertoni2011keccak_reference}. Las transformaciones $\theta$, $\rho$
y $\pi$ producen difusión y reorganización del estado, mientras que $\chi$
constituye la única capa no lineal de la ronda. La transformación $\iota$
incorpora una constante dependiente del número de ronda.

La capa $\chi$ opera localmente sobre grupos de cinco bits que comparten las
mismas coordenadas $(y,k)$. En este trabajo, cada uno de estos grupos se
interpreta como una S-box local de cinco bits. Una S-box se considera activa
cuando al menos uno de sus bits de entrada presenta una diferencia no nula
entre dos ejecuciones de Keccak.

Determinar manualmente el número mínimo de S-boxes activas resulta difícil
incluso para versiones reducidas de la permutación. Las diferencias son
reorganizadas por las transformaciones lineales y posteriormente atraviesan
la capa no lineal, lo que produce un espacio combinatorio que crece con el
tamaño del estado y con el número de rondas.

La Programación Lineal Entera Mixta
(\emph{Mixed-Integer Linear Programming}, MILP) permite representar este
problema mediante variables binarias, restricciones lineales y una función
objetivo. Esta metodología ha sido utilizada para automatizar la búsqueda de
características diferenciales y lineales en distintas primitivas
criptográficas \cite{mouha2012milp,sun2014automatic}.

El presente trabajo desarrolla un modelo MILP bit a bit para versiones
reducidas de Keccak con tamaños de palabra

\[
z \in \{4,8\}.
\]

Estas configuraciones producen estados de

\[
25z
\]

bits; por tanto, se estudian estados de 100 bits para $z=4$ y de 200 bits
para $z=8$. Debido a que existen cinco valores posibles de $y$ y $z$
posiciones dentro de cada palabra, cada ronda contiene

\[
5z
\]

aplicaciones locales de $\chi$. En consecuencia, los experimentos consideran
20 S-boxes por ronda para $z=4$ y 40 S-boxes por ronda para $z=8$.

El modelo representa simultáneamente dos ejecuciones concretas de Keccak.
Las diferencias se calculan mediante restricciones XOR exactas, en lugar de
utilizar una aproximación diferencial truncada de la capa $\chi$. Esta
decisión permite reconstruir parejas concretas de estados, verificar la
propagación bit a bit y validar independientemente los patrones de actividad
obtenidos.

La implementación se desarrolló en Python utilizando PuLP como biblioteca
de modelado y CBC como solver de programación entera. El análisis comprende
una, dos y tres rondas de Keccak reducido.


\subsection{Planteamiento del problema}
\label{subsec:planteamiento}

Sean $A_0^L$ y $A_0^R$ dos estados iniciales distintos de Keccak reducido.
Su diferencia inicial se define como

\[
\Delta A_0 = A_0^L \oplus A_0^R,
\]

con la condición

\[
\Delta A_0 \neq 0.
\]

Durante cada ronda, ambos estados atraviesan las transformaciones

\[
\theta \longrightarrow \rho \longrightarrow \pi
\longrightarrow \chi \longrightarrow \iota.
\]

Sea $\Delta B_r$ la diferencia entre las dos ejecuciones inmediatamente
antes de la capa $\chi$ en la ronda $r$. Una aplicación local de $\chi$,
identificada por las coordenadas $(y,k)$, se considera activa cuando

\[
\bigvee_{x=0}^{4} \Delta B_r[x,y,k] = 1.
\]

El problema consiste en minimizar la suma de S-boxes activas durante $R$
rondas, bajo la condición de que los estados iniciales sean diferentes.

La pregunta central del trabajo es la siguiente:

\begin{quote}
¿Cuál es el número mínimo de aplicaciones activas de la capa $\chi$ que
puede presentar una diferencia inicial no nula al propagarse durante una,
dos y tres rondas de Keccak reducido para $z=4$ y $z=8$?
\end{quote}

El número mínimo de S-boxes activas se denota por

\[
N_{\mathrm{act}}^{\min}(z,R),
\]

donde $z$ es el tamaño de palabra y $R$ es el número de rondas analizadas.

Para una y dos rondas se busca establecer mínimos globales mediante la
coincidencia entre una cota inferior demostrada y un testigo factible. Para
tres rondas, debido al crecimiento del espacio combinatorio, se busca
establecer cotas globales y construir testigos verificables mediante una
búsqueda exhaustiva restringida.


\subsection{Correspondencia experimental entre intentos y rondas}
\label{subsec:intentos_rondas}

La actividad académica propone relacionar tres rangos de intentos con
experimentos de una, dos y tres rondas. En este trabajo se adopta la
correspondencia

\[
R(t)=
\begin{cases}
1, & t < 10,\\
2, & 10 \leq t < 20,\\
3, & 20 \leq t < 30.
\end{cases}
\]

Esta relación se utiliza exclusivamente como criterio experimental para
seleccionar el número de rondas analizadas. El número de intentos no modifica
las transformaciones internas de Keccak y no representa una relación
criptográfica demostrada entre los intentos de un atacante y la profundidad
de la permutación.

Por esta razón, los resultados principales se expresan en términos del
tamaño de palabra $z$ y del número de rondas $R$.


\subsection{Objetivos}
\label{subsec:objetivos}

\subsubsection{Objetivo general}

Modelar una versión reducida de Keccak mediante Programación Lineal Entera
Mixta para determinar o acotar el número mínimo de S-boxes activas en una,
dos y tres rondas, utilizando CBC como solver.


\subsubsection{Objetivos específicos}

\begin{enumerate}
    \item Implementar las transformaciones $\theta$, $\rho$, $\pi$, $\chi$
    e $\iota$ para estados reducidos de Keccak con tamaños de palabra
    $z=4$ y $z=8$.

    \item Construir un modelo MILP bit a bit que represente dos ejecuciones
    simultáneas de Keccak y calcule exactamente sus diferencias mediante
    restricciones XOR.

    \item Definir variables binarias que identifiquen las aplicaciones
    activas de $\chi$ y formular una función objetivo que minimice su
    cantidad total.

    \item Analizar configuraciones de una, dos y tres rondas, distinguiendo
    entre mínimos globales, cotas inferiores, testigos factibles y óptimos
    restringidos.

    \item Validar los resultados mediante una implementación funcional
    independiente, reconstrucción de testigos y pruebas automatizadas.
\end{enumerate}


\subsection{Contribuciones}
\label{subsec:contribuciones}

Las principales contribuciones del trabajo son las siguientes:

\begin{enumerate}
    \item Una implementación funcional de las transformaciones
    $\theta$, $\rho$, $\pi$, $\chi$ e $\iota$ para versiones reducidas de
    Keccak.

    \item Una formulación MILP exacta que representa bit a bit dos
    ejecuciones de la permutación y calcula sus diferencias XOR.

    \item Un modelo de actividad que identifica cada aplicación local de
    $\chi$ activa y minimiza su suma durante varias rondas.

    \item La certificación de los mínimos globales

    \[
    N_{\mathrm{act}}^{\min}(4,1)
    =
    N_{\mathrm{act}}^{\min}(8,1)
    =
    1
    \]

    para una ronda, y

    \[
    N_{\mathrm{act}}^{\min}(4,2)
    =
    N_{\mathrm{act}}^{\min}(8,2)
    =
    4
    \]

    para dos rondas.

    \item La obtención de las cotas globales

    \[
    5
    \leq
    N_{\mathrm{act}}^{\min}(4,3)
    \leq
    13
    \]

    y

    \[
    5
    \leq
    N_{\mathrm{act}}^{\min}(8,3)
    \leq
    14
    \]

    para tres rondas.

    \item La construcción de testigos con distribuciones de actividad

    \[
    2+2+9
    \]

    para $z=4$ y

    \[
    2+2+10
    \]

    para $z=8$. Estos testigos son óptimos dentro de la familia restringida
    $2+2+c$ y proporcionan cotas superiores válidas para el problema
    global.

    \item Una estrategia de validación que contrasta las soluciones del
    modelo con una implementación funcional independiente y permite
    reconstruir los estados, las diferencias y los soportes activos.
\end{enumerate}


\subsection{Alcance}
\label{subsec:alcance}

El estudio se limita a versiones reducidas de Keccak con tamaños de palabra

\[
z \in \{4,8\}
\]

y a un máximo de tres rondas. Estas configuraciones conservan la estructura
de la permutación, pero no corresponden a la instancia completa
Keccak-$f[1600]$ utilizada por SHA-3.

La variable principal es el número de aplicaciones activas de la capa
$\chi$. Esta medida permite estudiar la propagación estructural de una
diferencia, pero no determina por sí sola la probabilidad de una
característica diferencial. El modelo no incorpora los pesos de las
transiciones diferenciales locales ni estima directamente la complejidad
de un ataque contra SHA-3.

Los resultados de una y dos rondas corresponden a mínimos globales
certificados. Para tres rondas, los valores 13 y 14 no se presentan como
mínimos globales. Estos valores son los mejores resultados obtenidos dentro
de la familia restringida

\[
2+2+c
\]

y deben interpretarse como cotas superiores globales respaldadas por
testigos concretos.

En consecuencia, el aporte del trabajo consiste en formular, resolver y
validar el problema de minimización de S-boxes activas para las
configuraciones reducidas consideradas, sin extender las conclusiones más
allá de las evidencias matemáticas y experimentales disponibles.